% ============================================================
% Appendix C: Annotation Guide
% Included from main.tex
% Sources: labeller/README.md; labeller/app.py (the rules enforced on
% save, MUTATIONS); Section~\ref{sec:custom-dataset}.
% ============================================================

\chapter{Annotation Guide}
\label{app:annotation-guide}

This is the full guide behind the custom validation set
(Section~\ref{sec:custom-dataset}), as the annotator and the AI
labelling pass apply it.

\section*{C.1\quad What is labelled}

One record holds one atomic, checkable claim, in the article's own
wording. A proposed sentence that asserts two things is split into two
records. A claim proposed by the daily batch is either labelled or
skipped with one of these reasons:

\begin{description}[leftmargin=2.6cm,labelwidth=2.4cm]
    \item[opinion] an opinion or interpretation, which no evidence can
    settle;
    \item[prediction] a statement about the future;
    \item[trivial] checkable but not worth checking;
    \item[fragment] not a whole sentence (a heading, a caption, a
    sentence cut short);
    \item[duplicate] the same claim as one already labelled;
    \item[other] none of the above, with a note.
\end{description}

\section*{C.2\quad Labels}

\begin{description}[leftmargin=2.6cm,labelwidth=2.4cm]
    \item[supported] (\texttt{TRUE}) the evidence supports the claim as
    stated.
    \item[partially supported] (\texttt{PARTIALLY\_TRUE}) the central
    claim holds and a detail does not.
    \item[conflicting] (\texttt{MISLEADING}) sources of equal standing
    conflict, or the claim is true but selected to mislead.
    \item[refuted] (\texttt{FALSE}) the evidence refutes the claim.
    \item[not enough evidence] (\texttt{UNVERIFIED}) there is not
    enough independent evidence to decide either way.
\end{description}

\section*{C.3\quad The four tie-break rules}

Fixed before labelling began; the note of a record names the rule that
decided a close call.

\begin{enumerate}[label=\textbf{Rule \arabic*.}, leftmargin=*]
    \item \textbf{The organisation's own source.} A claim whose only
    evidence is the communication of the organisation the story is about
    (a press release, report, website or post) is \emph{not enough
    evidence}, and the source is named in the note. \emph{Enforced on
    save.}
    \item \textbf{Numeric tolerance.} A figure within $\pm10\%$ of the
    source's, or a reasonable rounding of it, is supported. Beyond that,
    the claim is \emph{partially supported} if the figure is a detail of
    an otherwise correct claim, and \emph{refuted} if the figure is the
    claim.
    \item \textbf{Date of the evidence.} Only evidence available on the
    article's publication date is admissible. A claim confirmed only
    later stays \emph{not enough evidence}. \emph{Enforced on save}: an
    evidence date after the article's is refused.
    \item \textbf{Hierarchy of sources.} Primary source (report,
    official statistic, peer-reviewed study) $>$ reference media outlet
    $>$ press release $>$ social media. The higher tier prevails; sources
    of the same tier that conflict with none decisive make the claim
    \emph{conflicting}. The strongest tier used is recorded.
\end{enumerate}

Also enforced on save: any label other than \emph{not enough evidence}
needs at least one evidence link.

\section*{C.4\quad Procedure}

\begin{enumerate}
    \item Take the claim as proposed, or split it; keep the article's
    wording.
    \item Search for evidence independently of the system; the system's
    verdict is never shown.
    \item Record the evidence links, the date of the newest evidence used
    and the tier of the strongest.
    \item Apply the rules and choose the label; record the claimant, the
    claim type (factual, numerical or interpretive) and the topic.
    \item Explain any close call in the note, naming the rule.
\end{enumerate}

\section*{C.5\quad Synthetic refuted claims}

Only a natural claim labelled \emph{supported}, whose note states the
value the evidence gives, is mutated, and each source only once. One
element of what the claim asserts is substituted --- a number, a date,
an entity or a relation --- never a detail, and never by negation.
Figures that must agree are changed together. The substitute is
plausible without the evidence. Everything else is kept: the article,
date, topic, claim type, tier and evidence links, with the source's
identifier and the kind of substitution recorded. A mutation whose
source is later relabelled is listed for review.

\section*{C.6\quad AI-assisted labels}

A label proposed by the AI labelling pass starts its note with
``[AI label: Claude, \textit{date} --- not yet checked by the human
annotator]''. It follows this guide and cites the evidence it relied on.
It does not count as gold until the annotator has reviewed it and kept
or changed it; the agreement between the two is reported.

\section*{C.7\quad The blind review}

20\% of the records, the first in the order of the SHA-256 hash of
their identifiers, are labelled again at least a week after the first
pass, without the first label, the own-source flag or the note.
Agreement and Cohen's $\kappa$ are computed on the first label;
$\kappa \geq 0.6$ is the threshold for the guide to be considered
consistently applicable, and below it the rules behind the
disagreements are revised.
